\documentclass[11pt,fleqn]{book}

\input{../structure.tex}

\definecolor{mainColor}{RGB}{25, 102, 250}

\usepackage{booktabs}
\usepackage[table]{xcolor}
\usepackage[utf8]{inputenc}
\usepackage[italian]{babel}
\usepackage{amsmath}
\usepackage{amssymb}
% \usepackage{listings} % Per codice
% \usepackage{ragged2e}
\usepackage{graphicx}
% \usepackage{tikz}
\usetikzlibrary{calc, circuits.logic.US, positioning}
\tikzset{circuit logic US, every circuit symbol/.style={scale=1.2}}
\usepackage{booktabs}

\usepackage[none]{hyphenat}

\begin{document}

%----------------------------------------------------------------------------------------
%	TITLE PAGE
%----------------------------------------------------------------------------------------

\begingroup
\thispagestyle{empty}
\begin{tikzpicture}[remember picture,overlay]
    \coordinate [below=12cm] (midpoint) at (current page.north);
    \node at (current page.north west){
        \begin{tikzpicture}[remember picture,overlay]
            \node[anchor=north west,inner sep=0pt] at (0,0) {\includegraphics[width=\paperwidth]{libro.png}};
            \draw[anchor=north] (midpoint) node [fill=mainColor!30!white,fill opacity=0.9,text opacity=1,inner sep=1cm]{
                \Huge\centering\bfseries\sffamily\parbox[c][][t]{\paperwidth}
                {
                    \centering 
                    {
                        \Huge 
                        \ifsoluzioni
                            Soluzioni
                        \fi
                        Eserciziario
                    }
                    \\[15pt]
                    {\Huge TPS}
                    \\[15pt]
                    {\Large Classe III - informatica}
                    \\[20pt] 
                    {\huge Prof. Gallina Roberto}
                }
            };
        \end{tikzpicture}
    };
    \end{tikzpicture}
    \vfill
\endgroup

%----------------------------------------------------------------------------------------
%	COPYRIGHT PAGE
%----------------------------------------------------------------------------------------

\input{../copyright.tex}

%----------------------------------------------------------------------------------------
%	TABLE OF CONTENTS
%----------------------------------------------------------------------------------------

% \chapterimage{libro.jpg}

\pagestyle{empty}

\tableofcontents

\pagestyle{fancy}

\newpage

%----------------------------------------------------------------------------------------
%	PART 1
%----------------------------------------------------------------------------------------

% \partimage{chapter_images/sistema_binario.jpg}

\part{Codifica dell'informazione}

%----------------------------------------------------------------------------------------
%	CHAPTER 1
%----------------------------------------------------------------------------------------

% \chapterimage{chapter_head_1.pdf}

\chapter{Alfabeto sorgente e codice}

\newpage

\section{Codifica lunghezza fissa e variabile}

Dato il seguente alfabeto sorgente T e alfabeto codifica E,
\begin{itemize}
    \item determinare qual è la lunghezza minima per una codifica a lunghezza fissa
    \item fornire una possibile codifica a lunghezza fissa 
    \item fornire una possibile codifica a lunghezza variabile
    \item fornire una possibile codifica a lunghezza fissa raggruppando gli elementi in base a un criterio definito dallo studente
\end{itemize}
\emph{Non è necessario specificare tutte le combinazioni possibili, ma è necessario far capire come è costruita la codifica}
\vspace{0.5cm}

\begin{esercizio}[1]
    \textbf{Esercizio guidato}

    $T = \{ LUN, MAR, MER, GIO, VEN \} $
    
    $E = \{ G, H \} $ \\

    Poichè $|T| = 5$ e $|E| = 2$, applicando la formula posso calcolare la lunghezza minima per una codifica a lunghezza fissa:
        $m = \lceil log_{|E|}{(|T|)} \rceil = \lceil log_2{(5)} \rceil \approx 2.3 \xrightarrow{} 3$
        
    Infatti  $|E|^{m} = 2^3 = 8 > |T| = 5$ \\

    Posso proporre una possibile codifica a lunghezza fissa, in cui ogni codifica è lunghezza 3 caratteri dell'alfabeto codifica:
    \begin{itemize}
        \item LUN = GGG
        \item MAR = GGH
        \item MER = GHG
        \item GIO = GHH
        \item VEN = HGG
    \end{itemize}
    \vspace{0.25cm}

    Posso proporre una possibile codifica a lunghezza varibile, in cui ogni codifica è lunghezza massima di 3 caratteri dell'alfabeto codifica:
    \begin{itemize}
        \item LUN = G
        \item MAR = GH
        \item MER = HG
        \item GIO = HH
        \item VEN = H
    \end{itemize}
    \vspace{0.25cm}

    Posso proporre una possibile codifica a lunghezza fissa, dopo aver raggruppato gli elementi in base allo loro iniziale (M o altro). Ossia:

    \hspace{2cm}
    $T_M = \{ MAR, MER \} $
    \hspace{2cm}
    $T_{ALTRO} = \{ LUN, GIO, VEN \} $\\

    Avendo due gruppi avrà bisogno di $\lceil log_{|E|}{(numero\ gruppi)} \rceil = \lceil log_2{(2)} \rceil = 1$ cifre per rappresentare il gruppo; quindi:
    \begin{itemize}
        \item gli elementi del gruppo $T_M$ termineranno con $G$
        \item gli elementi del gruppo $T_{ALTRO}$ termineranno con $H$
    \end{itemize}
    \vspace{0.25cm}
    
    Poichè il gruppo più grande ha 3 elementi avrò bisogno di ulteriori  $\lceil log_2{(3)} \rceil \approx 1.58 \xrightarrow{} 2$ cifre per rappresentare gli elementi del gruppo:

    \begin{itemize}
        \item LUN = GGH
        \item MAR = GGG
        \item MER = GHG
        \item GIO = GHH
        \item VEN = HGH
    \end{itemize}

\end{esercizio}

\newpage

\eserciziopy{codifica_fissa_variabile}

% \section{Codifiche}

% \begin{esercizio}[1]
%     \textbf{Esercizio guidato}

%     Convertire il numero 80 dalla base 10 alla base 2

%     \lipsum[1]
% \end{esercizio}


%----------------------------------------------------------------------------------------
%	CHAPTER 2
%----------------------------------------------------------------------------------------

\chapter{Sistema binario}

\section{Base 2}

Convertire i seguenti numeri dalla base 10 alla base 2:

\begin{esercizio}[1]
    \textbf{Esercizio guidato}
    
    80

    \immediate\write18{python3 ./../scripts/Dec2Bin.py 80 > temp/ex_dec2bin80.tex}
    \input{./temp/ex_dec2bin80.tex}
\end{esercizio}

\eserciziopy{convDecBin}

\newpage

\noindent
Convertire i seguenti numeri dalla base 2 alla base 10:

\vspace{0.5cm}

\begin{esercizio}[1]
    \textbf{Esercizio guidato}

    10101

    \immediate\write18{python3 ./../scripts/Bin2Dec.py 10101 > temp/ex_bin2dec10101.tex}
    \input{./temp/ex_bin2dec10101.tex}
\end{esercizio}

\eserciziopy{convBinDec}

\newpage

\section{Base 2 - Modulo e Segno}

Convertire i seguenti numeri dalla base 10 alla base 2 usando la codifica modulo e segno:

\begin{esercizio}[1]
    \textbf{Esercizio guidato}

    Convertire il numero -71 dalla base 10 alla base 2

    \immediate\write18{python3 ./../scripts/Dec2BinMS.py "-71" > temp/ex_dec2bin2ms-71.tex}
    \input{./temp/ex_dec2bin2ms-71.tex}

\end{esercizio}

\eserciziopy{dec2bin_ms}

\newpage

\section{Base 2 - Complemento a 2}

Convertire i seguenti numeri dalla base 10 alla base 2 usando la codifica complemento a 2:

\begin{esercizio}[1]
    \textbf{Esercizio guidato}

    Convertire il numero -71 dalla base 10 alla base 2

    \immediate\write18{python3 ./../scripts/Dec2BinC2.py "-71" > temp/ex_dec2bin2c2-71.tex}
    \input{./temp/ex_dec2bin2c2-71.tex}

\end{esercizio}

\eserciziopy{dec2bin_c2}


% %----------------------------------------------------------------------------------------
% %	CHAPTER 3
% %----------------------------------------------------------------------------------------

% % \chapterimage{chapter_head_1.pdf}

% \chapter{Codifica parole}

% \section{ASCII}

% %----------------------------------------------------------------------------------------
% %	CHAPTER 4
% %----------------------------------------------------------------------------------------

% % \chapterimage{chapter_head_1.pdf}

% \chapter{Codifica immagini}

% \section{Convertire il colore RGB in Hex}

% \section{Convertire il colore Hex in RGB}

% \section{Calcolo profondità}

% %----------------------------------------------------------------------------------------
% %	CHAPTER 5
% %----------------------------------------------------------------------------------------

% % \chapterimage{chapter_head_1.pdf}

% \chapter{Codifica suoni}

% \section{Convertire il colore RGB in Hex}

% %----------------------------------------------------------------------------------------
% %	CHAPTER 6
% %----------------------------------------------------------------------------------------

% % \chapterimage{chapter_head_1.pdf}

% \chapter{Codifica video}

% \section{Convertire il colore RGB in Hex}

%----------------------------------------------------------------------------------------
%	PART 2
%----------------------------------------------------------------------------------------

% \partimage{chapter_images/sistema_binario.jpg}

\part{Operazioni sui dati}

%----------------------------------------------------------------------------------------
%	CHAPTER 7
%----------------------------------------------------------------------------------------

% \chapterimage{chapter_head_1.pdf}
\chapter{Operazioni binarie}

\section{Somma}

Effettua le seguenti somme in binario:

\begin{esercizio}[1]
    \textbf{Esercizio guidato}

    41 + 18

    Come prima cosa convertiamo i due numeri da decimale a binario

    \immediate\write18{python3 ./../scripts/Dec2Bin.py 41 > temp/ex_dec2bin41.tex}
    \input{./temp/ex_dec2bin41.tex}

    \immediate\write18{python3 ./../scripts/Dec2Bin.py 18 > temp/ex_dec2bin18.tex}
    \input{./temp/ex_dec2bin18.tex}

    A questo punto eseguire la somma bit a bit, ricordando che 1+1 fa 0 con riporto di 1 e 1+1+1 fa 1 con riporto di 1

    \immediate\write18{python3 ./../scripts/SommaBinaria.py "41+18" > temp/ex_somma.tex}
    \input{./temp/ex_somma.tex}

\end{esercizio}

\newpage

\eserciziopy{somme_binarie}

\section{Sottrazione}

Effettua le seguenti sottrazioni in binario:

\begin{esercizio}[1]
    \textbf{Esercizio guidato}

    41 - 18

    Come prima cosa convertiamo i due numeri da decimale a binario

    \immediate\write18{python3 ./../scripts/Dec2Bin.py 41 > temp/ex_dec2bin41.tex}
    \input{./temp/ex_dec2bin41.tex}

    \immediate\write18{python3 ./../scripts/Dec2Bin.py 18 > temp/ex_dec2bin18.tex}
    \input{./temp/ex_dec2bin18.tex}

    A questo punto eseguire la differenza bit a bit, ricordando che 0-1 fa 1 con il prestito di 1 dalla cifra precedente

    \immediate\write18{python3 ./../scripts/SottrazioneBinaria.py "41-18" > temp/ex_sottrazione.tex}
    \input{./temp/ex_sottrazione.tex}

\end{esercizio}

\newpage

\eserciziopy{sottrazioni_binarie}

\newpage

\section{Moltiplicazione}

Effettua le seguenti moltiplicazioni in binario:

\begin{esercizio}[1]
    \textbf{Esercizio guidato}

    41 x 9

    Come prima cosa convertiamo i due numeri da decimale a binario

    \immediate\write18{python3 ./../scripts/Dec2Bin.py 41 > temp/ex_dec2bin41.tex}
    \input{./temp/ex_dec2bin41.tex}

    \immediate\write18{python3 ./../scripts/Dec2Bin.py 9 > temp/ex_dec2bin9.tex}
    \input{./temp/ex_dec2bin9.tex}

    Andiamo ad eseguire la moltiplicazione del primo numero per ogni bit del secondo, ricordandosi di spostarsi di una posizione ogni volta.

    Al termine di deve eseguire la somma in colonna, qual'ora le righe risultimo molte conviene fare due righe per volta

    \immediate\write18{python3 ./../scripts/MoltiplicazioneBinaria.py "41x9" > temp/ex_moltiplicazione.tex}
    \input{./temp/ex_moltiplicazione.tex}

\end{esercizio}

\newpage

\eserciziopy{moltiplicazioni_binarie}

\newpage

\section{Divisione}

Effettua le seguenti divisioni in binario:

\begin{esercizio}[1]
    \textbf{Esercizio guidato}

    41 : 9

    Come prima cosa convertiamo i due numeri da decimale a binario

    \immediate\write18{python3 ./../scripts/Dec2Bin.py 41 > temp/ex_dec2bin41.tex}
    \input{./temp/ex_dec2bin41.tex}

    \immediate\write18{python3 ./../scripts/Dec2Bin.py 9 > temp/ex_dec2bin9.tex}
    \input{./temp/ex_dec2bin9.tex}

    Abbasso un bit alla volta fino a quando ottengo un numero maggiore o uguale del divisore, per ogni bit che abbasso mi segno uno 0 aggiunto un numero adatto mi segno un 1, scrivo il divisore sotto al numero stando attento di allineare le colonne e effettuo la sottrazione; ripeto le operazioni fino all'ultimo bit; 
    
    Il risultato sarà composto da quoziente e resto.

    \immediate\write18{python3 ./../scripts/DivisioneBinaria.py "41:9" > temp/ex_divisione.tex}
    \input{./temp/ex_divisione.tex}
\end{esercizio}

\newpage

\eserciziopy{divisioni_binarie}

%----------------------------------------------------------------------------------------
%	CHAPTER 8
%----------------------------------------------------------------------------------------

% \chapter{Rilevazione e correzione errori}

% \section{Bit di parità}

% \section{CRC}

% \section{LRC}

%----------------------------------------------------------------------------------------
%	PART 3
%----------------------------------------------------------------------------------------

% \partimage{chapter_images/sistema_binario.jpg}

\part{Sistema operativo}

%----------------------------------------------------------------------------------------
%	CHAPTER 9
%----------------------------------------------------------------------------------------

% \chapterimage{chapter_head_1.pdf}
\chapter{Scheduling}

\section{FCFS, SJF, SRTF}

Effettua lo schedeuling dei seguenti processi, usando l'algoritmo FCFS, SJF e SRTF:

\eserciziopy{sheduling1}

\newpage

\section{Round Robin}

Effettua lo schedeuling dei seguenti processi, usando l'algoritmo Round Robin con t=2 e t=3

\eserciziopy{sheduling2}

%----------------------------------------------------------------------------------------
%	CHAPTER 10
%----------------------------------------------------------------------------------------

% \chapterimage{chapter_head_1.pdf}
\chapter{Gestione memoria}

\section{Memoria a partizione fisse}

Mostra l'evoluzione della memoria dati i seguenti processi e la seguente memoria a partizioni fisse usando l'algoritmo FirstFit, BestFit e WorstFit

\eserciziopy{allocazione_fisse}

\newpage

\section{Memoria a partizione variabili}

Mostra l'evoluzione della memoria dati i seguenti processi e la seguente memoria a partizioni variabili usando l'algoritmo FirstFit, BestFit e WorstFit

\eserciziopy{allocazione_dinamiche}

\newpage

\section{Allocazione delle pagine}

Mostra l'evoluzione delle pagine allocate in memoria dati le seguenti pagine richieste usando gli algoritmi FIFO, LRU e OPT

\eserciziopy{allocazione_pagine}



\end{document}